\section{Basic functions of an echoscope}

    An echoscope is an ultrasonic measurement system that connects a probe to a computer or an oscilloscope. The probe is commonly made of a transducer in form 
    of a piezocristal. The echoscope then gives then an electical signal with a frequency to the probe which causes the piezocristal to vibrate and therefore
    it emits a supersonic wave. The probe not only emits ultrasound but it can act as a detector aswell when the piezocristal receives a vibration and converts
    it into an electical signal which then can be amplified and measured by the echoscope.\\
    There are two main meathods of measurements using a echoscope: Firstly one probe emits a wave and another probe receives the incoming wave. Here the two
    probes have to me aligned perfectly to not get disturbances. Secondly one probe emits a wave and the same probe also receives it using the reflection at
    the end surface. The probe has to be aligned perpendicular to the end surface to not get disturbances. \\
    There are multiple ways to process incoming signals from the echoscope. In this paper we want to highlight two. The most important one is the A-scan
    (amplitude scan) where the probe is kept at a constant position, measuring the amplitude of the emitted and received signal depending on the time after
    a pulse has been sent by the transmitter. Another way one can use an echoscope is the B-scan(brightness scan), where the probe is moved over an object
    and a 2D-Image of the plane in which teh probe was moved in is created. The amplitudes are given as different brightness values.